% GENERATED FILE. Edit canonical lessons, then run npm run generate:books.
% canonical-source-hash: fnv1a64:85623cdace8e9b48
% canonical-lessons: ES-C355-pared, ES-C355-techo, ES-C355-piso, ES-C355-rincon, ES-C355-fondo

\chapter{Inside The Building, Not Outside It}
\label{ch:la-pared-y-el-techo}
\hlchaptermodality{355}

\begin{chapteropening}
\textbf{By the end of this chapter:} \emph{I can say la pared, el techo, el piso, el rincon and el fondo, and place a thing inside a room instead of inside a town.}

Complete ES-C355-fondo: name the wall, the ceiling and the floor of a room, and put a thing in its corner or at its back.
\end{chapteropening}

\section[la pared]{la pared — the indoor wall, which Latin counted separately}

\label{lesson:ES-C355-pared}



\begin{quote}
Say \emph{the park}, then \emph{the garden}. (\emph{El parque}; \emph{el jardín}.)
\end{quote}

\subsection*{What to know first}
\begin{itemize}
\raggedright
  \item la casa — what it divides up.
  \item el cuarto — what it makes.
\end{itemize}

\begin{sounds}
\begin{itemize}
\raggedright
  \item \emph{pa-RED}, two beats with the weight on the last. The final \emph{d} is soft, nearly a breath. $\to$ pronunciation reference
\end{itemize}
\end{sounds}

\begin{cousinweb}
\textbf{la pared} — \textquotedblleft{}the wall,\textquotedblright{} the one inside a building.

Latin counted two kinds of wall and gave them unrelated names, and Spanish kept the pair.

\textbf{Pariēs} was a house wall — the thing that divides one room from the next, or holds a roof up. \textbf{Mūrus} was the wall around a town, the defensive one. A Roman would not have used one for the other.

Spanish still has both. \textbf{La pared} is the indoor wall you hang things on; \emph{el muro} and \emph{la muralla} are the ones around a city, and English borrowed the second family for \textbf{mural}, the painting made on a big wall.

\emph{Pariēs} had less luck in English, which kept it in one place only, and a strange one: the \textbf{parietal} bone is the one that walls the side of your skull. An anatomist's Latin, and the last survivor.
\end{cousinweb}

\begin{grammarlens}[title={What you've built}]
The first surface of a room, and a Latin distinction between indoors and out.
\end{grammarlens}

\subsection*{Your turn}
Say these aloud:

\begin{itemize}
\raggedright
  \item \emph{la pared}
  \item \emph{una pared}
  \item \emph{la pared de la casa}
\end{itemize}

\begin{tcolorbox}[breakable,colback=teal!4,colframe=teal!35!black,title={Before you move on}]
Which Latin wall was the one around a town? (\emph{Mūrus}.) And which one is \emph{la pared}? (The wall inside a house.)
\end{tcolorbox}
\section[el techo]{el techo — the covered part, and everything else that covers}

\label{lesson:ES-C355-techo}



\begin{quote}
Say \emph{the garden}, then \emph{the wall}. (\emph{El jardín}; \emph{la pared}.)
\end{quote}

\subsection*{What to know first}
\begin{itemize}
\raggedright
  \item la casa — what it covers.
  \item la nieve — what it keeps off you.
\end{itemize}

\begin{sounds}
\begin{itemize}
\raggedright
  \item \emph{TE-cho}, two beats with the weight on the first. The \emph{ch} is one sound, as in English \emph{cheese}. $\to$ pronunciation reference
\end{itemize}
\end{sounds}

\begin{cousinweb}
\textbf{el techo} — \textquotedblleft{}the roof,\textquotedblright{} and inside a room, \textquotedblleft{}the ceiling.\textquotedblright{}

Latin \textbf{tegere} is to cover, and \textbf{tēctum} is what has been covered — the covered part of a building. Spanish softened the middle of it into a \emph{ch} and got \emph{techo}.

The covering verb is one of the most productive Rome had, and English is full of it without noticing:

\begin{itemize}
\raggedright
  \item to \textbf{protect} is to cover in front of a thing;
  \item to \textbf{detect} is to take the cover off it;
  \item an \textbf{integument} is the covering of a body;
  \item a Roman \textbf{toga} was the thing he covered himself with;
  \item and a roof \textbf{tile} is a \textbf{tēgula}, a little coverer.
\end{itemize}

English has the root a second time, older and rougher. \textbf{Thatch} is the same ancient word arriving down the Germanic road instead of the Latin one, where the \emph{t} stayed put and the \emph{g} wore into a \emph{ch} — which is what a straw roof is: a covering.
\end{cousinweb}

\begin{grammarlens}[title={What you've built}]
The second surface of a room, and the verb behind half the covering words in English.
\end{grammarlens}

\subsection*{Your turn}
Say these aloud:

\begin{itemize}
\raggedright
  \item \emph{el techo}
  \item \emph{el techo de la casa}
  \item \emph{la pared y el techo}
\end{itemize}

\begin{tcolorbox}[breakable,colback=teal!4,colframe=teal!35!black,title={Before you move on}]
What does \emph{tegere} mean? (To cover.) And what does the \emph{de-} do in \emph{detect}? (Takes the cover off.)
\end{tcolorbox}
\section[el piso]{el piso — the floor is named for the treading}

\label{lesson:ES-C355-piso}



\begin{quote}
Say \emph{the wall}, then \emph{the ceiling}. (\emph{La pared}; \emph{el techo}.)
\end{quote}

\subsection*{What to know first}
\begin{itemize}
\raggedright
  \item el pie — what walks on it.
  \item la escoba — what sweeps it.
\end{itemize}

\begin{sounds}
\begin{itemize}
\raggedright
  \item \emph{PI-so}, two beats with the weight on the first. The \emph{s} is clear, never buzzed. $\to$ pronunciation reference
\end{itemize}
\end{sounds}

\begin{cousinweb}
\textbf{el piso} — \textquotedblleft{}the floor,\textquotedblright{} and in Spain, \textquotedblleft{}the flat.\textquotedblright{}

The floor is named after what you do to it. \emph{Pisar} is to tread, and it comes from Late Latin \textbf{pīnsāre}, to pound or crush — the verb for a mortar and pestle, not for a gentle step. The \textbf{piso} is the pounded part.

The pounding turns up in English twice, and both times in machinery:

\begin{itemize}
\raggedright
  \item a \textbf{pestle} is the thing that pounds in a mortar, from \textbf{pistillum}, the little pounder;
  \item a \textbf{piston} is a pounder in a cylinder, borrowed through Italian.
\end{itemize}

The word then did something practical. One trodden level of a building is \emph{un piso}, a storey — and in Spain a storey you can buy and live in is \emph{un piso}, a flat. Much of Latin America says \emph{el departamento} instead, and keeps \emph{piso} for the floor under your feet.

So the same three letters name the surface, the level and the home.
\end{cousinweb}

\begin{grammarlens}[title={What you've built}]
The third surface of a room, and one word doing the work of floor, storey and flat.
\end{grammarlens}

\subsection*{Your turn}
Say these aloud:

\begin{itemize}
\raggedright
  \item \emph{el piso}
  \item \emph{el piso de la casa}
  \item \emph{el techo y el piso}
\end{itemize}

\begin{tcolorbox}[breakable,colback=teal!4,colframe=teal!35!black,title={Before you move on}]
What does \emph{pisar} mean? (To tread.) And what does a pestle do? (It pounds.)
\end{tcolorbox}
\section[el rincón]{el rincón — the corner Arabic gave Spanish}

\label{lesson:ES-C355-rincon}



\begin{quote}
Say \emph{the ceiling}, then \emph{the floor}. (\emph{El techo}; \emph{el piso}.)
\end{quote}

\subsection*{What to know first}
\begin{itemize}
\raggedright
  \item la esquina — the corner out on the street.
  \item ojalá — another word Arabic gave Spanish.
\end{itemize}

\begin{sounds}
\begin{itemize}
\raggedright
  \item \emph{rin-KON}, two beats with the weight on the last, which is why it carries a written accent. The opening \emph{r} is rolled. $\to$ pronunciation reference
\end{itemize}
\end{sounds}

\begin{cousinweb}
\textbf{el rincón} — \textquotedblleft{}the corner,\textquotedblright{} the one inside a room.

Spanish keeps two corners apart, and it gets them from two different languages.

\textbf{La esquina}, which you have, is the corner you turn on a street — the outside angle of a building. It is Latin.

\textbf{El rincón} is the inside angle, where two \emph{paredes} meet and a chair goes. And it is Arabic: \textbf{rukn}, a corner, and also a supporting pillar. The corner of the Kaaba where the black stone sits is \emph{al-rukn}.

Eight centuries of al-Andalus left Spanish somewhere near four thousand Arabic words, and they are not exotic ones. You already own two: \textbf{ojalá}, which is \emph{wa-shā' allāh}, \textquotedblleft{}and may God will it,\textquotedblright{} and \textbf{hasta}, which is \emph{hattā}. Add \emph{el café}, from \emph{qahwa}, and \emph{el jabón} is in the same neighbourhood.

What is unusual about \emph{rincón} is that it displaced nothing and named something Latin had no separate word for. Arabic did not lend Spanish a synonym here. It lent it a distinction.
\end{cousinweb}

\begin{grammarlens}[title={What you've built}]
The corner indoors, and a reminder that a good part of Spanish did not come from Latin at all.
\end{grammarlens}

\subsection*{Your turn}
Say these aloud:

\begin{itemize}
\raggedright
  \item \emph{el rincón}
  \item \emph{un rincón}
  \item \emph{la silla está en el rincón}
\end{itemize}

\begin{tcolorbox}[breakable,colback=teal!4,colframe=teal!35!black,title={Before you move on}]
Which of the two corners is the Arabic one? (\emph{El rincón}, the one indoors.) And which Spanish word for hoping is also Arabic? (\emph{Ojalá}.)
\end{tcolorbox}
\section[el fondo]{el fondo — the bottom, and its Spanish twin}

\label{lesson:ES-C355-fondo}



\begin{quote}
Say \emph{the floor}, then \emph{the corner}. (\emph{El piso}; \emph{el rincón}.)
\end{quote}

\subsection*{What to know first}
\begin{itemize}
\raggedright
  \item dentro — where it is.
  \item f $\to$ h — the change that gave it a twin.
\end{itemize}

\begin{sounds}
\begin{itemize}
\raggedright
  \item \emph{FON-do}, two beats with the weight on the first. The \emph{d} after \emph{n} is firm, not the soft one between vowels. $\to$ pronunciation reference
\end{itemize}
\end{sounds}

\begin{cousinweb}
\textbf{el fondo} — \textquotedblleft{}the bottom,\textquotedblright{} and in a room, \textquotedblleft{}the back.\textquotedblright{}

Latin \textbf{fundus} is the bottom of a thing — of a jar, of a valley — and by extension the piece of ground a farm sits on. English keeps the bottom everywhere:

\begin{itemize}
\raggedright
  \item to \textbf{found} a building is to put its bottom in place, and that bottom is its \textbf{foundation};
  \item a \textbf{fund} is money held at the bottom of things, in reserve;
  \item a ship that \textbf{founders} goes to the bottom;
  \item and something \textbf{profound} is \emph{pro fundō}, toward the bottom.
\end{itemize}

Now the Spanish part. \textbf{Fondo} is a borrowed word — it came in from books, with its \emph{f} intact. But Latin \emph{fundus} also came down the long road, through mouths, and words that did that turned their Latin \emph{f} into a Spanish \emph{h}: \emph{facere} to \emph{hacer}, \emph{fabulari} to \emph{hablar}, \emph{ferrum} to \emph{hierro}.

Do it to \emph{fundus} and you get \textbf{hondo}, deep. So \emph{fondo} and \emph{hondo} are one Latin word that arrived twice, and the \emph{h} is the receipt for the longer journey.
\end{cousinweb}

\begin{grammarlens}[title={What you've built}]
A room named all the way round: \emph{la pared}, \emph{el techo}, \emph{el piso}, \emph{el rincón}, \emph{el fondo}.
\end{grammarlens}

\subsection*{Your turn}
Say these aloud:

\begin{itemize}
\raggedright
  \item \emph{el fondo}
  \item \emph{al fondo}
  \item \emph{la pared del fondo}
\end{itemize}

\begin{tcolorbox}[breakable,colback=teal!4,colframe=teal!35!black,title={Before you move on}]
What does \emph{fundus} mean? (The bottom.) And which Spanish word for \emph{deep} is the same word, arrived the long way? (\emph{Hondo}.)
\end{tcolorbox}
